% Auto-generated by build_latex.py — do not edit by hand.
% Test 2–tagged questions only

\subsection*{\texttt{test2\_sample\_questions\_q1}}
\textit{test2\_sample\_questions.pdf \#1 · test2}\par\smallskip
\begin{quote}
\small\sloppy
Jobs arrive at a system according to a Poisson process with rate 10 per hour. The service\\
    requirement of each job has an exponential distribution with mean 5 minutes. For use of\\
    the system, each arriving job is charged one dollar, with a 20 cent rebate if the job has\\
    to wait before entering service (the jobs wait in one queue). A server costs 1 dollar per\\
    hour to operate (this is a fixed cost, and does not depend on the server utilization). To\\
    maximize expected revenue for the system, should one or two servers be used?\\
\end{quote}

\subsection*{\texttt{test2\_sample\_questions\_q2}}
\textit{test2\_sample\_questions.pdf \#2 · test2}\par\smallskip
\begin{quote}
\small\sloppy
Short Answer Questions\\
\\
    (i) Very briefly explain why it is not a good idea to operate a queueing system under a\\
         heavy load (arrival rate close to service rate).\\
    (ii) Very briefly explain why an infinite buffer queue is said to be “unstable” when the\\
         mean arrival rate is greater than the mean service rate.\\
    (iii) A professor schedules student appointments every ten minutes. This is on the as-\\
          sumption that she can advise the average student in ten minutes. Will there be a\\
          long line outside of her office?\\
\end{quote}

\subsection*{\texttt{test2\_sample\_questions\_q3}}
\textit{test2\_sample\_questions.pdf \#3 · test2}\par\smallskip
\begin{quote}
\small\sloppy
A call centre is using the following staffing rule: the average time in queue (before service\\
    commences) of a customer must be less than 5 minutes. The minimum staff level should\\
    be kept. The time spent with a customer (service time) is exponentially distributed with\\
    mean 4 minutes and interarrival times are exponentially distributed.\\
    (a) For what range of arrival rates can one staff person be used?\\
    (b) Suppose the staffing rule was changed such that all customers should have a wait of\\
    less than 5 minutes. Is this achievable? Explain.\\
\end{quote}

\subsection*{\texttt{test2\_sample\_questions\_q4}}
\textit{test2\_sample\_questions.pdf \#4 · test2}\par\smallskip
\begin{quote}
\small\sloppy
Jobs arrive to a single server according to a Poisson process with rate 3 per hour. Service\\
    times are exponentially distributed with mean 1/µ hours.\\
    (a) Find µ so that the steady-state probability that there is more than 1 customer in the\\
    system is .25.\\
    (b) Now assume that µ = 5. The server gets a bonus of \$1 per job that does not have\\
    to wait in the queue before entering service. Calculate the expected bonus per hour that\\
    the server makes.\\
\end{quote}

\subsection*{\texttt{test2\_sample\_questions\_q5}}
\textit{test2\_sample\_questions.pdf \#5 · test2}\par\smallskip
\begin{quote}
\small\sloppy
Arrivals to a single server follow a Poisson process with rate 30 per hour. Service times\\
    are exponentially distributed with mean 60/35 minutes. What fraction of the arrival\\
    rate given in the question yields half the expected number in system? Explain why this\\
    fraction is more or less than one half.\\
\end{quote}

\subsection*{\texttt{test2\_sample\_questions\_q6}}
\textit{test2\_sample\_questions.pdf \#6 · test2}\par\smallskip
\begin{quote}
\small\sloppy
A system operates as follows: requests arrive according to a Poisson process with rate\\
    1 per second. There are two processors, each with exponentially distributed processing\\
    times with mean 1.25 seconds per request. The second processor only operates when\\
    there are at least T requests waiting for processing. The first processor always operates\\
\\
\\
\\
    if possible.\\
    (a) For T = 3, calculate the probability that an arriving request has to wait before being\\
    processed.\\
    (b) Is the value calculated in (a) an increasing or decreasing function of T ? Explain your\\
    answer.\\
\end{quote}

\subsection*{\texttt{test2\_sample\_questions\_q7}}
\textit{test2\_sample\_questions.pdf \#7 · test2}\par\smallskip
\begin{quote}
\small\sloppy
An existing system consists of a single server. Requests arrive according to a Poisson\\
    process with rate three per hour. The processing times are exponentially distributed with\\
    mean 18 minutes. You are asked whether it would be beneficial to add a second (identical)\\
    processor. The two processors would share a common queue. The costs are as follows:\\
    operating costs of 25 dollars per hour per processor and cost of delay of five dollars per\\
    hour per request.\\
\end{quote}

\subsection*{\texttt{test2\_sample\_questions\_q8}}
\textit{test2\_sample\_questions.pdf \#8 · test2}\par\smallskip
\begin{quote}
\small\sloppy
A system operates as follows. Arrivals occur according to a Poisson process with rate 1\\
    per minute. There is a single server and room for at most three jobs in the system. If\\
    there are one or two jobs in the system, the server has service times that are exponentially\\
    distributed with mean one minute. However, if there are three jobs in the system, the\\
    server speeds up so that the service times are exponentially distributed with mean 30\\
    seconds. Calculate the expected number of jobs in the system (in steady-state).\\
\end{quote}

\subsection*{\texttt{test2\_sample\_questions\_q9}}
\textit{test2\_sample\_questions.pdf \#9 · test2}\par\smallskip
\begin{quote}
\small\sloppy
The average waiting time (W ) on a database system is three seconds. During a one minute\\
    observation interval, the idle time on the system was measured to be 10 seconds. Using\\
    an M/M/1 model, determine the mean service time per query.\\
\end{quote}

\subsection*{\texttt{test2\_sample\_questions\_q10}}
\textit{test2\_sample\_questions.pdf \#10 · test2}\par\smallskip
\begin{quote}
\small\sloppy
A system works as follows. Arrivals occur according to a Poisson process with rate 2 per\\
    minute. If there are one or two jobs in the system, a single server processes the jobs,\\
    the processing times being exponentially distributed with mean 40 seconds. If there are\\
    three jobs in the system, two servers process the jobs, each server having processing times\\
    exponentially distributed with mean 40 seconds. Finally, an arriving job that sees three\\
    jobs in the system is rejected. In steady-state, calculate the expected number of jobs in\\
    the system.\\
\end{quote}

\subsection*{\texttt{test2\_sample\_questions\_q11}}
\textit{test2\_sample\_questions.pdf \#11 · test2}\par\smallskip
\begin{quote}
\small\sloppy
Arrivals occur to a system according to a Poisson process with rate one per minute. The\\
    blocking probability for the system is measured to be 0.1. The expected number of jobs\\
    in the system is 2.5. What is the expected response time?\\
\end{quote}

\subsection*{\texttt{test2\_sample\_questions\_q12}}
\textit{test2\_sample\_questions.pdf \#12 · test2}\par\smallskip
\begin{quote}
\small\sloppy
Students arrive at a single copy machine according to a Poisson process with rate 10\\
    students per hour. The time a student spends at the copying machine is exponentially\\
    distributed with a mean of 4 minutes.\\
    (a) What is the probability that a student arriving at the machine will have to wait?\\
    (b) What is the mean number of students waiting for the copier (not including the one\\
    making copies)?\\
\end{quote}

\subsection*{\texttt{test2\_sample\_questions\_q13}}
\textit{test2\_sample\_questions.pdf \#13 · test2}\par\smallskip
\begin{quote}
\small\sloppy
Consider a single processor implementing a priority discipline. There are two classes of\\
    jobs, high priority and low priority. Arriving jobs follow Poisson processes. The processing\\
    times are exponentially distributed.\\
\\
\\
\\
    (i) Low priority jobs are only processed if no high priority jobs are present. That is,\\
         any high priority jobs in the queue are served, one at a time, before the processor\\
         processes low priority jobs.\\
    (ii) If the queue only has low priority jobs, an arriving high priority job will instantly\\
         displace the low priority job at the processor and take its place.\\
\\
    Let the arrival rate be 5 jobs per millisecond for each class and the processing rate be 20\\
    jobs (of either class) per millisecond. Calculate the average number of high priority jobs\\
    in the system and the average number of jobs of both classes in the system. Use these two\\
    answers to deduce the average number of low priority jobs in the system. (Hint: Think\\
    carefully how high priority jobs are processed.)\\
\end{quote}

\subsection*{\texttt{test2\_sample\_questions\_q14}}
\textit{test2\_sample\_questions.pdf \#14 · test2}\par\smallskip
\begin{quote}
\small\sloppy
A single server systems has arrivals that occur according to a Poisson process with rate λ\\
    per hour. The processing times are exponentially distributed with mean 1/µ hours. The\\
    cost of waiting jobs is \$ C1 per job per hour. The cost of operating the system is \$ µC2\\
    per unit per hour, whether or not it is always in operation. How large should µ be to\\
    minimize the total cost?\\
\end{quote}

\subsection*{\texttt{test2\_sample\_questions\_q15}}
\textit{test2\_sample\_questions.pdf \#15 · test2}\par\smallskip
\begin{quote}
\small\sloppy
A designer is trying to decide how many processors to have in a system. Each processor\\
    works at a rate of 100 jobs per hour and the system is such that a processor upon becoming\\
    idle chooses the next waiting job. The system is being designed for an arrival rate of 150\\
    jobs per hour. Assume all underlying distributions (interarrival and service times) are\\
    exponential. If the design criterion is that the probability that an arriving job has to wait\\
    is no larger than .5, what is the smallest number of processors that can be chosen?\\
\end{quote}

\subsection*{\texttt{test2\_sample\_questions\_q16}}
\textit{test2\_sample\_questions.pdf \#16 · test2}\par\smallskip
\begin{quote}
\small\sloppy
Consider two M/M/1 systems with arrival rates λ1 = 0.5 per second and processing rates\\
    µ1 = 1 per second and an M/M/2 system with arrival rate λ2 = 2λ1 = 1 per second and\\
    processing rate µ2 = µ1 = 1 per second for each server. Compare the mean waiting time\\
    times for both systems. Discuss your result, giving a physical explanation for why one\\
    system is better.\\
\end{quote}

\subsection*{\texttt{test2\_sample\_questions\_q17}}
\textit{test2\_sample\_questions.pdf \#17 · test2}\par\smallskip
\begin{quote}
\small\sloppy
An M/M/1 system has a processing rate of 1 per second.\\
    (a) What is the maximum arrival rate that this system can support, if the requirement is\\
    that there are, on average, less than 4 jobs in the system, 90 percent of the time?\\
    (b) Suppose that at time t, there are 4 jobs in the system, with the job being processed\\
    having begun processing at time t − 0.5. What is the expected time of the next departure\\
    from the system?\\
\end{quote}

\subsection*{\texttt{test2\_sample\_questions\_q18}}
\textit{test2\_sample\_questions.pdf \#18 · test2}\par\smallskip
\begin{quote}
\small\sloppy
A system is described as follows. Arrivals occur according to a Poisson process with rate\\
    1 per second. There is room in the system for 3 jobs (including the one being processed).\\
    A single server has exponentially distributed processing times, but the rate depends on\\
    the number in system. A good fit is that when there are n jobs in the system, the\\
                      √\\
    processing rate is n jobs per second. Calculate the mean number of jobs in the system,\\
    in steady-state.\\
\end{quote}

\subsection*{\texttt{test2\_sample\_questions\_q19}}
\textit{test2\_sample\_questions.pdf \#19 · test2}\par\smallskip
\begin{quote}
\small\sloppy
(a) Consider an M/M/1/4 system. The observed mean number in system is 2, the mean\\
    waiting time (from arrival to departure) is 4 seconds. The arrival rate, λ (without taking\\
    the losses into account) is λ = 0.75 jobs per second. What is the loss (blocking) proba-\\
    bility?\\
    (b) Consider an M/M/c system with arrival rate 5 per minute. Each processor has a mean\\
    processing time of 20 seconds. How many processors are required so that the system is\\
    stable (i.e. the steady-state probabilities pn exist)?\\
    (c) If you could convert the system in (b) to c M/M/1’s, where an arrival joins each queue\\
    with probability 1/c, would you make the change? Briefly explain why or why not.\\
\end{quote}

\subsection*{\texttt{test2\_sample\_questions\_q20}}
\textit{test2\_sample\_questions.pdf \#20 · test2}\par\smallskip
\begin{quote}
\small\sloppy
Consider two M/M/1 systems with arrival rates λ1 = 0.5 per second and processing rates\\
    µ1 = 1 per second and an M/M/2 system with arrival rate λ2 = 2λ1 = 1 per second and\\
    processing rate µ2 = µ1 = 1 per second for each server. Compare the mean waiting time\\
    times for both systems. Discuss your result, giving a physical explanation for why one\\
    system is better.\\
\end{quote}

\subsection*{\texttt{test2\_sample\_questions\_q21}}
\textit{test2\_sample\_questions.pdf \#21 · test2}\par\smallskip
\begin{quote}
\small\sloppy
You are asked to design a system for which you receive 1 dollar in revenue for every\\
    arriving job that is served without waiting before processing and 50 cents for every job\\
    that is processed but has to wait before processing. Arrivals occur according to a Poisson\\
    process with rate 54 per hour and the processing times are exponentially distributed with\\
    mean 1 minute. You are considering the following two choices:\\
\\
     (a) A single processor, admit all arrivals, processing in FIFO order.\\
    (b) A single processor, limit the number in system to 5, serve arrivals in FIFO order.\\
\\
    Which would you choose if your goal is to maximize the rate of revenue?\\
\end{quote}

\subsection*{\texttt{test2\_sample\_questions\_q22}}
\textit{test2\_sample\_questions.pdf \#22 · test2}\par\smallskip
\begin{quote}
\small\sloppy
A single processor has room for at most three in the system. Processing times are exponen-\\
    tially distributed with mean 1, while the interarrival times are modelled as exponentially\\
    distributed with the following rates:\\
\\
                                 number in system arrival rate\\
                                        0             1.5\\
                                        1             1.0\\
                                        2             0.5\\
\\
    Calculate the utilization of the processor.\\
\end{quote}

\subsection*{\texttt{test2\_sample\_questions\_q23}}
\textit{test2\_sample\_questions.pdf \#23 · test2}\par\smallskip
\begin{quote}
\small\sloppy
There are two independent arrival streams (following independent Poisson processes) with\\
    respective arrival rates λ1 = 20 and λ2 = 15 per hour. Processing times for both types\\
    of jobs are exponentially distributed with mean 2 minutes. You have two processors and\\
    have the choice of dedicating one processor to each arrival stream, or combining the two\\
    arrival streams into a single queue and serving the arrivals in a FIFO manner at the first\\
    available processor of the two. Which configuration would you choose? Quantify the\\
    resulting performance difference between the two configurations.\\
\end{quote}

\subsection*{\texttt{test2\_sample\_questions\_q24}}
\textit{test2\_sample\_questions.pdf \#24 · test2}\par\smallskip
\begin{quote}
\small\sloppy
An M/M/1/3 system has an arrival rate of 8 per hour and mean processing time of 6\\
    minutes. The profit received from each arrival that is admitted into the system is 5\\
    dollars. For an extra 60 dollars per day, you can convert it to an M/M/1/4. Would you\\
    make the change?\\
\end{quote}

\subsection*{\texttt{test2\_sample\_questions\_q25}}
\textit{test2\_sample\_questions.pdf \#25 · test2}\par\smallskip
\begin{quote}
\small\sloppy
As a first cut at dimensioning a database server, we assume that the server has arrivals\\
    occurring as a Poisson process with rate 2 per minute and that the processing times are\\
    exponentially distributed with rate µ. If one pays a fixed cost of µ per minute for the\\
    capital investment and 10 for each arrival that has to wait to begin processing, what\\
    would you choose for µ?\\
\end{quote}

\subsection*{\texttt{test2\_sample\_questions\_q26}}
\textit{test2\_sample\_questions.pdf \#26 · test2}\par\smallskip
\begin{quote}
\small\sloppy
A system has room for at most two in system (i.e. arrivals finding two in the system are\\
    lost). Arrivals follow a Poisson process with rate 1 per minute. The processing times are\\
    exponentially distributed with mean 40 seconds if there is one job in the system and 20\\
    seconds if there are two jobs in the system. Find the expected number of lost arrivals in\\
    a one hour interval.\\
\end{quote}

\subsection*{\texttt{test2\_sample\_questions\_q27}}
\textit{test2\_sample\_questions.pdf \#27 · test2 · T/F}\par\smallskip
\begin{quote}
\small\sloppy
True or False.\\
    (a) One M/M/4 queue has lower mean waiting time in system than 4 M/M/1 queues.\\
    (The arrival rate is 4λ to the M/M/4, while the arrival rate is λ to each M/M/1. All\\
    servers have rate µ.)\\
    (b) One M/M/4 queue has lower average server utilization than 4 M/M/1 queues. (The\\
    arrival rate is 4λ to the M/M/4, while the arrival rate is λ to each M/M/1. All servers\\
    have rate µ.)\\
    (c) A system has an arrival rate of 2 per minute and an average waiting time in system\\
    of 10 minutes. The average number in system is 20.\\
\end{quote}

\subsection*{\texttt{test2\_sample\_questions\_q28}}
\textit{test2\_sample\_questions.pdf \#28 · test2}\par\smallskip
\begin{quote}
\small\sloppy
You are setting up a server system that has a single queue. Service is FCFS and you need\\
    to decide the number of servers. For cost reasons, you would like to choose the smallest\\
    number of servers. Arrivals follow a Poisson process with rate 20 per hour. Processing\\
    times are exponentially distributed with mean 2 minutes. There is a requirement that\\
    there is at most one job in the system 75 percent of the time. How many servers would\\
    you choose?\\
\end{quote}

\subsection*{\texttt{test2\_sample\_questions\_q29}}
\textit{test2\_sample\_questions.pdf \#29 · test2}\par\smallskip
\begin{quote}
\small\sloppy
A system is modelled as follows. There is space for at most two jobs in the system.\\
    Arrivals occur according to a Poisson process with rate 1 per minute. Arrivals finding\\
    two in the system are lost. If there is one job in the system, it is processed at rate 2\\
    per minute. If there are two in the system, due to contention, they are each processed\\
    at rate 0.75 per minute. Each job has its own processing time and processing times are\\
    exponentially distributed. Calculate the steady-state probability that there is exactly one\\
    job in the system.\\
\end{quote}

\subsection*{\texttt{test2\_sample\_questions\_q30}}
\textit{test2\_sample\_questions.pdf \#30 · test2}\par\smallskip
\begin{quote}
\small\sloppy
An administrator is trying to calculate how many (identical) servers are required for\\
    a database server system. Suppose arrivals occur according to a Poisson process with\\
    rate 1 per second. The administrator estimates that the processing time of a request is\\
\\
\\
\\
    exponentially distributed with mean 0.8 seconds. If the SLA (Service Level Agreement)\\
    is that the mean response time be less than 3 seconds, how many servers are required?\\
\end{quote}

\subsection*{\texttt{test2\_sample\_questions\_q31}}
\textit{test2\_sample\_questions.pdf \#31 · test2}\par\smallskip
\begin{quote}
\small\sloppy
A system operates as follows. Arrivals occur to a single server according to a Poisson\\
    process with rate 2 per minute. If there is one job in the system, its processing time is\\
    exponentially distributed with mean 30 seconds. If there are two jobs in the system, job\\
    processing times are exponentially distributed with mean 20 seconds. If an arriving job\\
    finds two jobs in the system, it is lost.\\
    (a) Calculate the steady-state probability that the server is idle.\\
    (b) Suppose that the last arrival occurred twenty seconds ago. How long would you expect\\
    to wait until the next arrival occurs?\\
\end{quote}

\subsection*{\texttt{test2\_sample\_questions\_q32}}
\textit{test2\_sample\_questions.pdf \#32 · test2}\par\smallskip
\begin{quote}
\small\sloppy
The performance of a system with fault detection is evaluated as follows. When operating,\\
    faults occur at a rate of λ = 1 per hour, with fault times being exponentially distributed.\\
    With probability 0.99, a fault is detected and recovery takes an exponentially distributed\\
    period of time with mean 1 second. If the fault is not detected, recovery takes an ex-\\
    ponentially distributed period of time with mean 5 minutes. Calculate the steady-state\\
    probability that the system is operating.\\
\end{quote}

\subsection*{\texttt{test2\_sample\_questions\_q33}}
\textit{test2\_sample\_questions.pdf \#33 · test2}\par\smallskip
\begin{quote}
\small\sloppy
Consider a system with arrivals following a Poisson process and having processing times\\
    that are exponentially distributed. Suppose that the state is the number of jobs in the\\
    system and the only non-zero rates are:\\
\\
                                     λ0 = λ\\
                                           λ\\
                                     λn =    ,         n = 1, 2\\
                                          n2\\
                                     µn = µ,         n = 1, 2, 3.\\
\\
    Calculate the steady-state probability that the system has less than 2 jobs present if λ = 2\\
    per second and µ = 1 per second.\\
\end{quote}

\subsection*{\texttt{test2\_sample\_questions\_q34}}
\textit{test2\_sample\_questions.pdf \#34 · test2}\par\smallskip
\begin{quote}
\small\sloppy
A system operates with a central server that processes incoming requests. The operating\\
    cost of the server is 20 dollars per hour, whether it is working or not. The cost of a\\
    waiting user is 15 dollars per hour per user. If the processing time of a job is exponentially\\
    distributed with mean three minutes and the arriving requests follow a Poisson process\\
    with rate 15 per hour, should another server be purchased?\\
\end{quote}

\subsection*{\texttt{test2\_sample\_questions\_q35}}
\textit{test2\_sample\_questions.pdf \#35 · test2}\par\smallskip
\begin{quote}
\small\sloppy
A system is described as follows. Arrivals occur according to a Poisson process with rate\\
    10 per minute. There is room in the system for 4 jobs (including the one being processed).\\
    A single server has exponentially distributed processing times, but the rate increases with\\
    the number in system. When there are n jobs in the system, the processing rate is 4n jobs\\
    per minute. Calculate the utilization of the server and the probability that an arriving\\
    job cannot enter the system.\\
\end{quote}

\subsection*{\texttt{test2\_sample\_questions\_q36}}
\textit{test2\_sample\_questions.pdf \#36 · test2 · T/F}\par\smallskip
\begin{quote}
\small\sloppy
True or False.\\
\\
\\
\\
     (a) One M/M/10 queue has lower mean number in system than 10 M/M/1 queues.\\
         (The arrival rate is 10λ to the M/M/10 system and λ to each of the M/M/1’s. The\\
         processing rate is µ for each server.)\\
\end{quote}

\subsection*{\texttt{test2\_sample\_questions\_q37}}
\textit{test2\_sample\_questions.pdf \#37 · test2}\par\smallskip
\begin{quote}
\small\sloppy
A service system consists of a single server, where arrivals to the server occur according to\\
    a Poisson process with rate three per hour. Processing times are exponentially distributed\\
    with mean 7/24 hours. The operating costs of the server are \$7.50 per hour and the cost\\
    due to waiting jobs is 50 cents per hour per job waiting. A design change could reduce\\
    the mean processing time by five minutes, but operating costs would be increased. How\\
    much of an increase in operating costs could be tolerated?\\
\end{quote}

\subsection*{\texttt{test2\_sample\_questions\_q38}}
\textit{test2\_sample\_questions.pdf \#38 · test2}\par\smallskip
\begin{quote}
\small\sloppy
Consider a system with two identical processors. There are two independent arrival\\
    streams following Poisson processes with rate λ1 = 25 per hour to the first processor and\\
    λ2 = 10 per hour to the second processor. Each processor has mean processing time of\\
    two minutes and processing times are exponentially distributed.\\
    (a) Calculate the mean response time at each processor.\\
    (b) Suppose that the arrival streams are combined and we have one queue for both\\
    processors, does the mean response time decrease for all arrivals? Justify your answer\\
    with appropriate calculations.\\
\end{quote}

\subsection*{\texttt{test2\_sample\_questions\_q39}}
\textit{test2\_sample\_questions.pdf \#39 · test2}\par\smallskip
\begin{quote}
\small\sloppy
A system operates with a central server that processes incoming requests. The operating\\
    cost of a server is 20 dollars per hour. whether it is working or not. The cost of waiting\\
    requests is 15 dollars per hour per request. If the processing times are exponentially\\
    distributed with mean three minutes and the arriving requests follow a Poisson process\\
    with rate 15 per hour, should another server be purchased?\\
\end{quote}

\subsection*{\texttt{test2\_sample\_questions\_q40}}
\textit{test2\_sample\_questions.pdf \#40 · test2}\par\smallskip
\begin{quote}
\small\sloppy
A web server is modelled as follows. We assume that it is a single server where arrivals\\
    follow a Poisson process with rate λ and the processing times are exponentially distributed\\
    with rate µ. The mean response time of the server is 10 seconds and its utilization is 0.8.\\
    What are λ and µ?\\
\end{quote}

\subsection*{\texttt{test2\_sample\_questions\_q41}}
\textit{test2\_sample\_questions.pdf \#41 · test2}\par\smallskip
\begin{quote}
\small\sloppy
Consider the following system. Arrivals occur according to a Poisson process with rate\\
    1 per minute if there are less than 2 jobs in the system and at rate 1 every 2 minutes if\\
    there are 2 jobs in the system. If there are 3 jobs in the system, arrivals are rejected.\\
    There are 3 identical servers - processing times are exponentially distributed with mean\\
    30 seconds. Calculate the steady-state probability that 2 or more servers are idle.\\
\end{quote}

\subsection*{\texttt{test2\_sample\_questions\_q42}}
\textit{test2\_sample\_questions.pdf \#42 · test2}\par\smallskip
\begin{quote}
\small\sloppy
Arrivals occur to a system of two servers (each with its own queue), according to a Poisson\\
    process with rate 1 per minute. The processing times for each server are exponentially\\
    distributed with mean 90 seconds. With probability 1/2, an arrival is sent to each of the\\
    servers.\\
\\
     (a) The last arrival to the first server occurred 1 minute ago. What is the expected time\\
         until the next arrival to the first server?\\
     (b) Calculate the mean response time at each queue.\\
\end{quote}

\subsection*{\texttt{test2\_sample\_questions\_q43}}
\textit{test2\_sample\_questions.pdf \#43 · test2}\par\smallskip
\begin{quote}
\small\sloppy
As a first cut at dimensioning a database server, we assume that the srever has arivals\\
    according to a Poisson process with rate 2 per minute and that the processing times are\\
    exponentially distributed with rate µ. If one pays a fixed cost of µ per minute for the\\
    cost of the server and a cost of 10 for each arrival that has to wait to begin processing,\\
    what would you choose for µ?\\
\end{quote}

\subsection*{\texttt{test2\_sample\_questions\_q44}}
\textit{test2\_sample\_questions.pdf \#44 · test2}\par\smallskip
\begin{quote}
\small\sloppy
(a) A remote sensing station has three identical sensors. The three sensors fail inde-\\
        pendently, where the failure times are exponentially distributed with mean 90 days.\\
        When all three sensors are failed, a repairman is sent, and the time to repair all\\
        three sensors (simultaneously) is exponentially distributed with mean 3 days. What\\
        is the steady-state probability that all three sensors are operating?\\
     (b) What is the rate (in steady-state) of visits of the repairman to the station?\\
\end{quote}

\subsection*{\texttt{test2\_sample\_questions\_q45}}
\textit{test2\_sample\_questions.pdf \#45 · test2}\par\smallskip
\begin{quote}
\small\sloppy
Suppose we have a system that has two processors in parallel with a single queue. The\\
    arrivals follow a Poisson process with rate 100 per hour and each processor has processing\\
    times that are exponentially distributed with rate 75 per hour. We would like to replace\\
    the two processors with a single processor. What processing rate is required for the single\\
    processor so that the mean response time is unchanged (assume that the processing times\\
    remain exponentially distributed)?\\
\end{quote}

\subsection*{\texttt{test2\_sample\_questions\_q46}}
\textit{test2\_sample\_questions.pdf \#46 · test2}\par\smallskip
\begin{quote}
\small\sloppy
A simplified game server is modelled as follows. Players arrive to the server according to\\
    a Poisson process with rate of one every 2 minutes. Arrivals wait until three players have\\
    arrived, at which point the game immediately starts. The game takes an exponentially\\
    distributed period of time with mean 20 minutes. When the game is complete, all of\\
    the players leave the system. This infinitely repeats, i.e. another game is started after\\
    three more players arrive. Potential players that arrive when a game is in progress leave\\
    without waiting.\\
\\
     (a) Calculate the limiting probability that a game is being played.\\
     (b) Suppose that we add room so that one player can wait for the next game while the\\
         current game is in progress (three players are still required for a game). Draw the\\
         transition rate diagram (do not solve for the limiting probabilities).\\
\end{quote}

\subsection*{\texttt{test2\_sample\_questions\_q47}}
\textit{test2\_sample\_questions.pdf \#47 · test2}\par\smallskip
\begin{quote}
\small\sloppy
You are asked to design a system for which you receive 1 dollar in revenue for every\\
    arriving job that is served without waiting before processing and 50 cents for every job\\
    that is processed but has to wait before processing. Arrivals occur according to a Poisson\\
    process with rate 54 per hour and the processing times are exponentially distributed with\\
    mean 1 minute. You are considering the following two choices:\\
\\
     (a) A single processor, admit all arrivals, processing in FIFO order.\\
     (b) A single processor, limit the number in system to 5, serve arrivals in FIFO order.\\
\\
    Which would you choose if your goal is to maximize the rate of revenue?\\
\end{quote}

\subsection*{\texttt{test2\_sample\_questions\_q48}}
\textit{test2\_sample\_questions.pdf \#48 · test2 · T/F}\par\smallskip
\begin{quote}
\small\sloppy
True or False.\\
\\
\\
\\
     (a) One needs λ < µ in an M/M/∞ queue for the limiting distribution to exist.\\
     (b) An M/M/1/4 queue can be modelled as a Birth-Death process.\\
     (c) When the number of servers in an M/M/k system becomes large, k = 2λ/µ servers\\
         is a good choice to get reasonable performance while keeping k as low as possible.\\
     (d) In an M/M/1 queue, for some choice of arrival rate, it is possible to double the\\
         expected response time with a one percent increase in the arrival rate.\\
\end{quote}

\subsection*{\texttt{test2\_sample\_questions\_q49}}
\textit{test2\_sample\_questions.pdf \#49 · test2}\par\smallskip
\begin{quote}
\small\sloppy
Consider an M/M/2 system with arrival rate 2 per second and processing times with mean\\
    1.2 seconds. If we replace this by two M/M/1 queues with the overall arrival rate of 2 per\\
    second and random routing (an arriving job is equally likely to go to either queue), what\\
    would the processing rates of the M/M/1 queues need to be so that the mean response\\
    time for a job remains unchanged?\\
\end{quote}

\subsection*{\texttt{test2\_sample\_questions\_q50}}
\textit{test2\_sample\_questions.pdf \#50 · test2}\par\smallskip
\begin{quote}
\small\sloppy
A computer server experiences one failure per 1000 hours of operation. Upon failure\\
    detection, a diagnostic process is run. In 95 percent of the cases, the cause of failure can\\
    be removed based on the diagnostic and service is restored in (on average) 15 minutes. In\\
    the remaining 5 percent of the cases, a repair person must be summoned and the expected\\
    time to service restoration is 12 hours. If you assume that all underlying distributions are\\
    exponential, calculate the probability that the server is working.\\
\end{quote}

\subsection*{\texttt{test2\_sample\_questions\_q51}}
\textit{test2\_sample\_questions.pdf \#51 · test2}\par\smallskip
\begin{quote}
\small\sloppy
A simple game server operates as follows. Two players are required to play the game,\\
    and only one game can be played at a time. Players arrive to the server following a\\
    Poisson process at rate 4 per hour. If a player arrives and finds a game in progress, that\\
    player does not enter the system. Once two players are waiting, the game immediately\\
    commences. The time to play a game is exponentially distributed with mean 30 minutes\\
    and both players exit the system once the game completes. Furthermore, players are\\
    impatient. That is, when a single player is waiting for another to arrive and the game to\\
    start, if they wait longer than an exponentially distributed random variable with mean\\
    30 minutes, they exit the system. The server is idle when no game is being played. What\\
    is the steady-state probability that the server is idle?\\
\end{quote}

\subsection*{\texttt{test2\_sample\_questions\_q52}}
\textit{test2\_sample\_questions.pdf \#52 · test2}\par\smallskip
\begin{quote}
\small\sloppy
Consider an M/M/2 system with λ = 1.5 and µ = 1. We wish to replace this with an\\
    M/M/1 system (with the same λ). We want to keep the probability that a job waits before\\
    beginning processing to be the same for both systems. What must be the processing rate\\
    of the M/M/1 system?\\
\end{quote}

\subsection*{\texttt{test2\_sample\_questions\_q53}}
\textit{test2\_sample\_questions.pdf \#53 · test2}\par\smallskip
\begin{quote}
\small\sloppy
A firm attracts new clients according to a Poisson process with rate one per week. Each\\
    client remains with the firm for an exponentially distributed period of time with mean 40\\
    weeks. In steady-state, what is the expected number of clients that the firm has?\\
\end{quote}

\subsection*{\texttt{test2\_sample\_questions\_q54}}
\textit{test2\_sample\_questions.pdf \#54 · test2}\par\smallskip
\begin{quote}
\small\sloppy
You are asked to determine the number of servers that a system should have. Requests\\
    arrive to the system according to a Poisson process with rate two per minute. The mean\\
    time to process a request on a server is exponentially distributed with mean 20 seconds.\\
\\
\\
\\
    If an arriving request finds all of the servers busy, it is lost. Determine the minimum\\
    number of servers needed such that the probability of losing a request is at most 0.15.\\
\end{quote}

\subsection*{\texttt{test2\_sample\_questions\_q55}}
\textit{test2\_sample\_questions.pdf \#55 · test2}\par\smallskip
\begin{quote}
\small\sloppy
Consider a system with two processors. Arrivals to the system occur according to a\\
    Poisson process with rate 2. The first processor works at rate 2, the second processor\\
    at rate 1 (processing times are exponentially distributed). The system is a loss system,\\
    so if both processors are busy, arrivals are lost. Also, if the system is empty, the faster\\
    processor is assigned to an arrival. Once a processor starts processing a job, it cannot be\\
    transferred to the other processor. Determine the steady-state probability that the slower\\
    processor is busy.\\
\end{quote}

\subsection*{\texttt{test2\_sample\_questions\_q56}}
\textit{test2\_sample\_questions.pdf \#56 · test2}\par\smallskip
\begin{quote}
\small\sloppy
(a) Consider an M/M/1 system with arrival rate 10. If the performance requirement is\\
        that there are no more than two jobs in the system 95 percent of the time, what is\\
        the required processing rate, µ.\\
     (b) For your choice of µ in (a), would an M/M/∞ model be a good approximation of the\\
         M/M/1 model? Here, a good approximation is if the mean response time is within\\
         10 percent of the actual value. Justify your answer.\\
\end{quote}

\subsection*{\texttt{test2\_sample\_questions\_q57}}
\textit{test2\_sample\_questions.pdf \#57 · test2}\par\smallskip
\begin{quote}
\small\sloppy
Suppose that we have a system with three machines, each with failure rate 1 per day.\\
    Repair times have mean 2 hours, and repairs are made by a single repair person. Suppose\\
    that it costs 400 dollars per day for a repair person (whether they are busy or not) and\\
    10,000 dollars per day whenever the number of operating machines drops below two. (For\\
    example, if half of the time less than two machines are operating, then you pay 5,000\\
    dollars per day.) Should a second repair person be added?\\
\end{quote}

\subsection*{\texttt{test2\_sample\_questions\_q58}}
\textit{test2\_sample\_questions.pdf \#58 · test2}\par\smallskip
\begin{quote}
\small\sloppy
Consider a single server that has 2 GB of memory. Two kinds of arrivals occur to the\\
    server: type 1 arrivals that require 1 GB of memory and type 2 arrivals that require 2\\
    GB of memory. Type 1 arrivals follow a Poisson process with rate 1 per minute and Type\\
    2 arrivals follow a Poisson process with rate 0.5 per minute. Processing times for type 1\\
    jobs are exponentially distributed with mean 30 seconds. Processing times for type 2 jobs\\
    are exponentially distributed with mean 1 minute. Any arrivals that need more memory\\
    than is available are rejected. Calculate the utilization of the server and the probability\\
    that a type 1 arrival is rejected.\\
\end{quote}

\subsection*{\texttt{test2\_sample\_questions\_q59}}
\textit{test2\_sample\_questions.pdf \#59 · test2}\par\smallskip
\begin{quote}
\small\sloppy
Consider a system with two users and one server. The two users have think times that are\\
    exponentially distributed with rates λ1 and λ2 , respectively. The CPU uses FCFS and\\
    has processing times that are exponentially distributed with rate µ. Draw an appropriate\\
    transition rate diagram for this system.\\
\end{quote}

\subsection*{\texttt{test2\_sample\_questions\_q60}}
\textit{test2\_sample\_questions.pdf \#60 · test2}\par\smallskip
\begin{quote}
\small\sloppy
A system is described by an M/M/1 queue with rates λ = 1 and µ = 2. You are considering\\
    a design change that would result in a self-service system (each job is its own server), but\\
    the mean processing time would increase. What is the maximum mean processing time\\
    allowable so that there is an improvement in the mean response time?\\
\end{quote}

\subsection*{\texttt{test2\_sample\_questions\_q61}}
\textit{test2\_sample\_questions.pdf \#61 · test2}\par\smallskip
\begin{quote}
\small\sloppy
A system operates with a central server that processes requests from a large number of\\
    users. The operating cost of the server is 20 dollars per hour, whether it is working or\\
    not. The cost of a waiting user is 15 dollars per hour per user.\\
    (a) If the processing time of a job is exponentially distributed with mean three minutes\\
    and the arriving requests follow a Poisson process with rate 15 per hour, should another\\
    server be purchased?\\
    (b) If the processing times and interarrival times are both deterministic (the values are\\
    the same as the means in (a)), would your recommendation change?\\
\end{quote}

\subsection*{\texttt{test2\_sample\_questions\_q62}}
\textit{test2\_sample\_questions.pdf \#62 · test2}\par\smallskip
\begin{quote}
\small\sloppy
Consider a system with three nodes. Arrivals occur to node 1 according to a Poisson\\
    process of rate 2 per second. Arrivals occur to node 2 according to a Poisson process\\
    of rate 1 per second. After processing at node 1, jobs return to node 1 with probability\\
    0.25, otherwise they go to node 3. After processing at node 2, jobs return to node 2 with\\
    probability 0.2, otherwise they go to node 3. After processing at node 3, jobs leave the\\
    system. The service policy at each node is FCFS. The processing times at each node\\
    are exponentially distributed with rates 3 per second, 2 per second and 5 per second,\\
    respectively.\\
    (a) Calculate the expected number of jobs at node 3.\\
    (b) You wish to improve the expected time in system for jobs that originally arrive to\\
    node 2. To do this, you are allowed to increase the processing rate of one node by ten\\
    percent. Which processor would you improve?\\
\end{quote}

\subsection*{\texttt{test2\_sample\_questions\_q63}}
\textit{test2\_sample\_questions.pdf \#63 · test2}\par\smallskip
\begin{quote}
\small\sloppy
We consider a system that allows a constant level of N programs (jobs) sharing main\\
    memory. It has 2 servers, server 1 being the CPU and server 2 being the I/O. The mean\\
    service time for a job at the CPU is 0.02 seconds, the mean service time for a job at the\\
    I/O subsystem is 0.04 seconds (assume all times are exponentially distributed). After a\\
    job completes at the I/O subsystem, it always goes to the CPU. After completion of a\\
    job at the CPU, a job exits the system with probability 0.05, otherwise it goes to I/O.\\
    When a job exits a system, another immediately enters at the CPU, thus there are always\\
    exactly N jobs in the system. For N = 2, calculate the utilization of the CPU (proportion\\
    of time CPU is busy).\\
\end{quote}

\subsection*{\texttt{test2\_sample\_questions\_q64}}
\textit{test2\_sample\_questions.pdf \#64 · test2}\par\smallskip
\begin{quote}
\small\sloppy
A network consists of three single server nodes. External arrivals to the first node follow\\
    a Poisson process with rate 10 per minute. After service at node 1, jobs go to node 2 with\\
    probability .5 and go to node 3 otherwise. After service at node 2, jobs return to node\\
    2 with probability .7 and go to node 3 otherwise. After service at node 3, jobs always\\
    exit the system. The service times at each node are exponentially distributed, with the\\
    following means: Node 1 - 5 seconds, Node 2 - 2 seconds, Node 3 - 4 seconds.\\
    (a) Find the expected number of customers at node 2.\\
    (b) If you could increase the service rate at one node, which node would be your last\\
    choice to do so?\\
\end{quote}

\subsection*{\texttt{test2\_sample\_questions\_q65}}
\textit{test2\_sample\_questions.pdf \#65 · test2}\par\smallskip
\begin{quote}
\small\sloppy
A closed network operates as follows. It has 4 nodes, the processing rates being 5, 10,\\
    10, 1 at the nodes, respectively. After finishing processing at node 1, jobs proceed to\\
\\
\\
\\
    either node 2 or 3 with equal probability. After processing at either nodes 2 or 3, with\\
    probability 1/10 a job visits node 4, otherwise it returns to node 1. After processing at\\
    node 4, jobs return to node 1.\\
    Calculate the throughput of this system with 2 jobs circulating. If you were to improve one\\
    processing rate, which would it be? Would you expect your recommendation to change\\
    as we increase the number of jobs in the system?\\
\end{quote}

\subsection*{\texttt{test2\_sample\_questions\_q66}}
\textit{test2\_sample\_questions.pdf \#66 · test2}\par\smallskip
\begin{quote}
\small\sloppy
Consider a closed queueing network with N = 3 nodes and M = 2 jobs. The processing\\
    times are exponentially distributed with rates µ1 = 0.72, µ2 = 0.64, µ3 = 1.00. The\\
    routing probabilities are r31 = 0.4, r32 = 0.6, r13 = 1.0, r23 = 1.0.\\
    (a) Calculate the throughput of the system.\\
    (b) Which node is the bottleneck?\\
\end{quote}

\subsection*{\texttt{test2\_sample\_questions\_q67}}
\textit{test2\_sample\_questions.pdf \#67 · test2}\par\smallskip
\begin{quote}
\small\sloppy
Construct a closed queueing network with two jobs and two servers such that the through-\\
    put is one per second. (You are free to choose all network parameters other than the\\
    number of jobs and servers.)\\
\end{quote}

\subsection*{\texttt{test2\_sample\_questions\_q68}}
\textit{test2\_sample\_questions.pdf \#68 · test2}\par\smallskip
\begin{quote}
\small\sloppy
A queueing network is described as follows. Arrivals occur to node one according to a\\
    Poisson process with rate one per minute. Processing times at node one are exponentially\\
    distributed with mean 20 seconds. After exiting node one, jobs go to nodes two or three\\
    with equal probability. The processing times at nodes two and three are exponentially\\
    distributed with means of 60 and 70 seconds, respectively. After processing at node two\\
    or three, a job leaves the system with probability 0.75, otherwise it returns to node one.\\
    Calculate the network bottleneck.\\
\end{quote}

\subsection*{\texttt{test2\_sample\_questions\_q69}}
\textit{test2\_sample\_questions.pdf \#69 · test2}\par\smallskip
\begin{quote}
\small\sloppy
Consider a network where external arrivals occur to node 1 at a rate of γ jobs per second.\\
    After processing at node 1, a job is routed to node 2 with probability 0.4 or to node 3\\
    with probability 0.3. Otherwise, it leaves the network. After processing at node 2 or node\\
    3, jobs always return to node 1. Processing times at nodes 1, 2, and 3 are exponentially\\
    distributed with rates 10 per second, 2 per second, and 2 per second, respectively.\\
    (a) If the utilization of node 3 is measured to be 0.7, what is γ?\\
    (b) Using the value of γ from (a), if the bottleneck processor has its processing rate\\
    doubled, calculate the mean number of jobs in the network.\\
\end{quote}

\subsection*{\texttt{test2\_sample\_questions\_q70}}
\textit{test2\_sample\_questions.pdf \#70 · test2}\par\smallskip
\begin{quote}
\small\sloppy
A system has four single server FCFS nodes, the CPU (node one), Printer (node two), Disk\\
    (node three), and an Input Device (node four). Processing times at each of the nodes\\
    are exponentially distributed with means 1/µ1 = 0.04 seconds, 1/µ2 = 0.03 seconds,\\
    1/µ3 = 0.06 seconds, and 1/µ4 = 0.05 seconds. Arrivals to the system occur to node 4\\
    at a rate of 4 jobs per second. The non-zero routing probabilities are p12 = p13 = 0.5,\\
    p41 = p21 = 1, and p31 = 0.6.\\
    (a) Compute the probability that there are exactly three jobs at the CPU, two jobs at\\
    the printer, four jobs at the disk, and one job at the input device.\\
    (b) Calculate the mean time that a job spends in the system (from arrival to departure).\\
\\
\\
\\
    (c) If the bottleneck device has its processing rate doubled, is there a new bottleneck\\
    device? If so, which device is it?\\
\end{quote}

\subsection*{\texttt{test2\_sample\_questions\_q71}}
\textit{test2\_sample\_questions.pdf \#71 · test2}\par\smallskip
\begin{quote}
\small\sloppy
Consider the following closed queueing network. There are three nodes, each having a\\
    single processor and exponentially distributed processing times. Departures from node 1\\
    are equally likely to go to the other two nodes. Departures from node 2 go to node 1 with\\
    probability 0.8, otherwise they go to node 3. Departures from node 3 always go to node\\
    1. The mean processing times at the three nodes are 0.5, 1 and 1.2, respectively. Two\\
    jobs circulate in the network.\\
    (a) What is the system bottleneck?\\
    (b) Calculate the throughput at node 1.\\
    (c) Calculate the mean waiting time at node 2.\\
\end{quote}

\subsection*{\texttt{test2\_sample\_questions\_q72}}
\textit{test2\_sample\_questions.pdf \#72 · test2}\par\smallskip
\begin{quote}
\small\sloppy
Consider the following network. Arrivals occur from outside to node 1 according to a\\
    Poisson process with rate 2 per minute. After processing at node 1, jobs go to node 2\\
    with probability 1/3 and to node 3 with probability 2/3. After processing at node 2, jobs\\
    return to node 1 with probability 1/2, otherwise they leave the network. After processing\\
    at node 3, jobs return to node 1 with probability 1/2, otherwise they leave the network.\\
    All processing times are exponentially distributed, each node has a single processor that\\
    uses FCFS. The mean processing times are node 1 - 7.5 seconds, node 2 - 40 seconds and\\
    node 3 - 15 seconds.\\
    (a) Find the probability that there is exactly one job in the network.\\
    (b) If you were to improve the processing rate at exactly one node, which would it be?\\
\end{quote}

\subsection*{\texttt{test2\_sample\_questions\_q73}}
\textit{test2\_sample\_questions.pdf \#73 · test2}\par\smallskip
\begin{quote}
\small\sloppy
Consider a closed queueing network with N = 3 nodes and M = 2 jobs. Processing\\
    times at the three nodes are exponentially distributed with rates µ1 = 0.7, µ2 = 0.1 and\\
    µ3 = 1.5. After processing at nodes 1 or 2, jobs always go back to node 3. After processing\\
    at node 3, jobs go to node 1 with probability 0.4 and node 2 with probability 0.6.\\
    (a) Calculate the bottleneck node.\\
    (b) Calculate the throughput at node 3.\\
\end{quote}

\subsection*{\texttt{test2\_sample\_questions\_q74}}
\textit{test2\_sample\_questions.pdf \#74 · test2}\par\smallskip
\begin{quote}
\small\sloppy
Consider a closed queueing network with N = 3 nodes and M = 2 jobs. Processing\\
    times at the three nodes are exponentially distributed with rates µ1 = 0.7, µ2 = 0.1 and\\
    µ3 = 1.5. After processing at nodes 1 or 2, jobs always go back to node 3. After processing\\
    at node 3, jobs go to node 1 with probability 0.4 and node 2 with probability 0.6.\\
    (a) Calculate the bottleneck node.\\
    (b) Calculate the throughput at node 2.\\
\end{quote}

\subsection*{\texttt{test2\_sample\_questions\_q75}}
\textit{test2\_sample\_questions.pdf \#75 · test2}\par\smallskip
\begin{quote}
\small\sloppy
Consider the following open network. External arrivals occur to node 1 according to a\\
    Poisson process with rate 2 per second. External arrivals occur to node 2 according to\\
    a Poisson process with rate 1 per second. After processing at node 1, jobs immediately\\
    return to node 1 with probability 0.25, otherwise they go to node 3. After processing\\
    at node 2, jobs immediately return to node 2 with probability 0.2, otherwise they go to\\
    node 3. After processing at node 3, jobs leave the system. The processing times at each\\
\\
\\
\\
    node are exponentially distributed with rate µ per second, 2 per second and 5 per second,\\
    respectively. There is one processor at each node.\\
    (a) For what range of values of µ is node 1 the bottleneck (while still keeping the system\\
    stable)?\\
    (b) Let the processing rate at node 1 be twice the maximum value in (a). What is the\\
    steady-state probability that all three processors are simultaneously busy?\\
\end{quote}

\subsection*{\texttt{test2\_sample\_questions\_q76}}
\textit{test2\_sample\_questions.pdf \#76 · test2}\par\smallskip
\begin{quote}
\small\sloppy
A closed network operates as follows. It has 4 nodes, the processing rates being 5, 10,\\
    10, 1 at the nodes, respectively. After finishing processing at node 1, jobs proceed to\\
    either node 2 or 3 with equal probability. After processing at either nodes 2 or 3, with\\
    probability 1/10 a job visits node 4, otherwise it returns to node 1. After processing at\\
    node 4, jobs return to node 1.\\
    Calculate the throughput of this system with 2 jobs circulating. If you were to improve one\\
    processing rate, which would it be? Would you expect your recommendation to change\\
    as we increase the number of jobs in the system?\\
\end{quote}

\subsection*{\texttt{test2\_sample\_questions\_q77}}
\textit{test2\_sample\_questions.pdf \#77 · test2}\par\smallskip
\begin{quote}
\small\sloppy
We consider a system that allows a constant level of N programs (jobs) sharing main\\
    memory. It has 2 servers, server 1 being the CPU and server 2 being the I/O. The mean\\
    service time for a job at the CPU is 0.02 seconds, the mean service time for a job at the\\
    I/O subsystem is 0.04 seconds (assume all times are exponentially distributed). After a\\
    job completes at the I/O subsystem, it always goes to the CPU. After completion of a\\
    job at the CPU, a job exits the system with probability 0.05, otherwise it goes to I/O.\\
    When a job exits a system, another immediately enters at the CPU, thus there are always\\
    exactly N jobs in the system. For N = 2, calculate the utilization of the CPU (proportion\\
    of time CPU is busy).\\
\end{quote}

\subsection*{\texttt{test2\_sample\_questions\_q78}}
\textit{test2\_sample\_questions.pdf \#78 · test2}\par\smallskip
\begin{quote}
\small\sloppy
An open queueing network is described as follows. External arrivals occur to node 1\\
    according to a Poisson process with rate 1 per minute and to node 2 according to a\\
    Poisson process with rate 3 per minute. After completing processing at node 1, a job\\
    is equally likely to go to nodes 2 or 3. After completing processing at node 2, a job is\\
    equally likely to go to nodes 1 or 3. After completing processing at node 3, a job goes to\\
    node 2 with probability 1/2, otherwise it leaves the network. Each node has processing\\
    rate 12 per minute, with processing times being exponentially distributed.\\
    (a) Calculate the steady-state probability that there is exactly one job in the network.\\
    (b) Calculate the mean number of jobs at node 1.\\
\end{quote}

\subsection*{\texttt{test2\_sample\_questions\_q79}}
\textit{test2\_sample\_questions.pdf \#79 · test2}\par\smallskip
\begin{quote}
\small\sloppy
(a) A closed network has two jobs circulating. There are three queues, with a single\\
        server at each. The first server has processing rate µ per second, the second and\\
        third servers both have processing rates of 1 per second. All processing times are\\
        exponentially distributed. The routing probabilities are P12 = 0.4, P13 = 0.6 and\\
        P21 = P31 = 1. For what range of µ is server 1 the bottleneck?\\
    (b) Choose a µ such that server 1 is the bottleneck. For your choice of µ, what is the\\
        probability that server 2 is idle?\\
\end{quote}

\subsection*{\texttt{test2\_sample\_questions\_q80}}
\textit{test2\_sample\_questions.pdf \#80 · test2}\par\smallskip
\begin{quote}
\small\sloppy
Consider a closed network with two nodes and M jobs circulating. The processing rate at\\
    the first node is 1 job per minute, while for the second node it is 3 jobs per minute. The\\
    network is cyclic, i.e. the routing probabilities are P12 = P21 = 1. All processing times\\
    are exponentially distributed. What is the minimum value of M such that the first node\\
    has a utilization of at least 0.9? Justify your answer.\\
\end{quote}

\subsection*{\texttt{test2\_sample\_questions\_q81}}
\textit{test2\_sample\_questions.pdf \#81 · test2}\par\smallskip
\begin{quote}
\small\sloppy
Consider a network with three nodes. Arrivals occur to the first node according to a\\
    Poisson process with rate four per hour. When jobs are finished processing at the first\\
    node, they always go to the second node. When finished processing at the second node,\\
    jobs exit the network with probability 0.25, go to the third node with probability 0.25,\\
    or return to the second node with probability 0.5. After processing at the third node,\\
    jobs always exit the network. Processing times are exponentially distributed, with the\\
    following rates: six per hour at the first node, 16 per hour at the second node, and 2.5\\
    per hour at the third node.\\
\\
     (a) What is the expected amount of time that a job spends in the network (from arrival\\
         to exiting)?\\
     (b) If the bottleneck node had its processing rate doubled, how would your answer in\\
         part (a) change?\\
\end{quote}

\subsection*{\texttt{test2\_sample\_questions\_q82}}
\textit{test2\_sample\_questions.pdf \#82 · test2}\par\smallskip
\begin{quote}
\small\sloppy
A network consists of three single server nodes. External arrivals to the first node follow\\
    a Poisson process with rate 10 per minute. After service at node 1, jobs go to node 2 with\\
    probability .5 and go to node 3 otherwise. After service at node 2, jobs return to node\\
    2 with probability .7 and go to node 3 otherwise. After service at node 3, jobs always\\
    exit the system. The service times at each node are exponentially distributed, with the\\
    following means: Node 1 - 5 seconds, Node 2 - 2 seconds, Node 3 - 4 seconds.\\
\\
     (a) Find the expected number of jobs at node 2.\\
     (b) If you could increase the service rate at one node, which node would be your last\\
         choice to do so?\\
\end{quote}

\subsection*{\texttt{test2\_sample\_questions\_q83}}
\textit{test2\_sample\_questions.pdf \#83 · test2}\par\smallskip
\begin{quote}
\small\sloppy
Consider a system of three single-server, infinite-buffer nodes. From the outside, arrivals\\
    occur to node 1 according to a Poisson process with rate 5 per minute. Jobs departing\\
    node 1 are equally likely to go to node 2 or node 3. Jobs departing node 2 or node 3 are\\
    equally likely to exit the network or return to node 1. Processing times at the nodes are\\
    exponentially distributed with mean 5 seconds (node 1), 8 seconds (node 2), and 9 seconds\\
    (node 3). Determine which node is the bottleneck and also calculate the probability that\\
    the network is empty.\\
\end{quote}

\subsection*{\texttt{test2\_sample\_questions\_q84}}
\textit{test2\_sample\_questions.pdf \#84 · test2}\par\smallskip
\begin{quote}
\small\sloppy
Consider a web server modelled as an open network. The server has one CPU and two\\
    disks. Arrivals occur to the CPU according to a Poisson process with rate 2 per second.\\
    After being processed at the CPU, with probability 0.02 a job leaves the network, with\\
    probability 0.40 it goes to disk A, and with probability 0.58 it goes to disk B. After being\\
    processed at either of the disks, a job goes to the CPU. The processing times (per visit)\\
\\
\\
\\
    at the CPU are exponentially distributed with rate 200 per second. At the two disks, the\\
    processing times (per visit) are exponentially distributed with rate 60 per second.\\
    (a) If you were to improve the processing rate of one device, which would it be?\\
    (b) Suppose that you balanced the loads on the disks. What would be the resulting mean\\
    number of jobs in the system?\\
\end{quote}

\subsection*{\texttt{test2\_sample\_questions\_q85}}
\textit{test2\_sample\_questions.pdf \#85 · test2}\par\smallskip
\begin{quote}
\small\sloppy
Consider a closed queueing network with three nodes and two jobs circulating. The\\
    routing probabilities are P12 = 1.0, P21 = P23 = P31 = P32 = 0.5. The processing rates\\
    are µ1 = µ2 = µ3 = 1. Processing times at all nodes are exponentially distributed. Find\\
    the expected number of jobs at the bottleneck node.\\
\end{quote}

\subsection*{\texttt{test2\_sample\_questions\_q86}}
\textit{test2\_sample\_questions.pdf \#86 · test2}\par\smallskip
\begin{quote}
\small\sloppy
A server provides downloads for two files, file A and file B. Requests follow a Poisson\\
    process with rate 0.15 per second. A request is for file A with probability 0.75 and for\\
    file B with probability 0.25. The time for a file A download is exponentially distributed\\
    with mean 8 seconds and the time for a file B download is exponentially distributed with\\
    mean 3 seconds. A request that arrives to the server when it is performing a download is\\
    lost. Find the probability (in steady-state) that the server is busy.\\
\end{quote}

\subsection*{\texttt{test2\_sample\_questions\_q87}}
\textit{test2\_sample\_questions.pdf \#87 · test2}\par\smallskip
\begin{quote}
\small\sloppy
Arrivals occur to a common queue (of infinite size) with rate 0.7 per minute. Processing\\
    times are exponentially distributed with mean one minute.\\
     (a) If the requirement is that the probability that an arriving job waits to begin pro-\\
         cessing is less than 0.2, what is the minimum number of servers that can be used?\\
     (b) Suppose that budget constraints are such that only one server can be used, so we\\
         decide to use a finite buffer instead. What size buffer would you choose if you wish\\
         to maximize the throughput while ensuring that that the probability that a job waits\\
         to begin processing is less than 0.2?\\
\end{quote}

\subsection*{\texttt{test2\_sample\_questions\_q88}}
\textit{test2\_sample\_questions.pdf \#88 · test2}\par\smallskip
\begin{quote}
\small\sloppy
A closed network has 3 nodes and 2 jobs circulating. After visiting node 1, a job is equally\\
    likely to visit nodes 2 or 3. After processing at node 2 or node 3, jobs always return to\\
    node 1. The processing times at nodes 1, 2 and 3 are exponentially distributed with\\
    means 1, 1, and 3, respectively. Calculate the probability that the bottleneck node is idle.\\
\end{quote}

\subsection*{\texttt{test2\_sample\_questions\_q89}}
\textit{test2\_sample\_questions.pdf \#89 · test2}\par\smallskip
\begin{quote}
\small\sloppy
Consider the following system. There are two servers in parallel, each can hold exactly one\\
    job (the job that it is processing). The processing times at each server are exponentially\\
    distributed with mean 30 seconds. Arrivals to the system occur as follows.\\
     (i) Type 1 arrivals follow a Poisson process with rate 2 per minute. If Server 1 is idle,\\
         they are assigned to Server 1, otherwise they are rejected.\\
     (ii) Type 2 arrivals follow a Poisson process with rate 2 per minute. If both servers are\\
          idle, the job is replicated on both servers (with independent processing times). If\\
          only one server is idle, it is assigned to that server. If both servers are busy, the job\\
          is rejected.\\
    (iii) If both servers are processing a replicated Type 2 job, and one server completes, the\\
          remaining replica is immediately killed.\\
\\
\\
\\
    Calculate the throughput of both Type 1 and Type 2 jobs.\\
\end{quote}

\subsection*{\texttt{test2\_sample\_questions\_q90}}
\textit{test2\_sample\_questions.pdf \#90 · test2}\par\smallskip
\begin{quote}
\small\sloppy
Jobs arrive to node one of a network according to a Poisson process with rate 1.25 per\\
    minute. With probability 0.4, jobs exiting node one go to node two, with probability\\
    0.5, jobs exiting node one go to node three, otherwise they exit the system. Jobs exiting\\
    node two go to node three or node one with equal probability. Jobs exiting node three\\
    go to node two or node one with equal probability. Processing times at nodes one, two\\
    and three are exponentially distributed with means 3 seconds, 5 seconds and 4 seconds,\\
    respectively. Determine\\
\\
     (a) the network bottleneck\\
    (b) the mean time that a job spends in the network, from arrival to departure\\
\end{quote}

\subsection*{\texttt{test2\_sample\_questions\_q91}}
\textit{test2\_sample\_questions.pdf \#91 · test2}\par\smallskip
\begin{quote}
\small\sloppy
A job’s execution is modelled as follows. To complete the job, three tasks must be\\
    completed. Tasks 1 and 2 can be done in parallel, requiring times that are exponentially\\
    distributed with rates µ1 and µ2 , respectively. After Tasks 1 and 2 are both complete,\\
    Task 3 is executed, requiring a time that is exponentially distributed with rate µ3 . A\\
    system continuously executes these jobs, so that when one job is completed, another\\
    immediately begins.\\
\\
     (a) Show how a CTMC model can be used to determine the steady-state probability\\
         that Task 3 is executing. (Construct the CTMC and give the equations that must\\
         be solved to determine the required probability.)\\
    (b) In terms of your answer for (a), what would the throughput of the system be (in\\
        jobs per time unit)?\\
\end{quote}

\subsection*{\texttt{test2\_sample\_questions\_q92}}
\textit{test2\_sample\_questions.pdf \#92 · test2}\par\smallskip
\begin{quote}
\small\sloppy
Suppose that we have an M/M/2 system where both servers are busy and there is one\\
    job waiting in the queue. What is the probability that the job waiting in the queue leaves\\
    the system before one of the jobs currently being processed?\\
\end{quote}

\subsection*{\texttt{test2\_sample\_questions\_q93}}
\textit{test2\_sample\_questions.pdf \#93 · test2}\par\smallskip
\begin{quote}
\small\sloppy
A system is operating with a single queue and two servers. The arrival rate is λ = 1 and\\
    the processing rate of each of the servers is µ = 2/3. All underlying distributions are\\
    exponential. We wish to replace the two servers by a single server working at rate µ∗ .\\
    What must µ∗ be to ensure that the probability that a job waits before being processed\\
    is equal to that of the two-server system?\\
\end{quote}

\subsection*{\texttt{test2\_sample\_questions\_q94}}
\textit{test2\_sample\_questions.pdf \#94 · test2}\par\smallskip
\begin{quote}
\small\sloppy
Consider the following network of single-server nodes. Arrivals occur from outside to\\
    node 1 according to a Poisson process with rate 1 job per second. The processing times\\
    at node 1 are exponentially distributed with rate 3 jobs per second. Jobs departing node\\
    1 always go to node 2. Node 2 processing times are exponentially distributed with rate 5\\
    jobs per second. After departing node 2, with probability 0.25 a job departs the network,\\
    otherwise it goes to node 3. Processing times at node 3 are exponentially distributed with\\
    rate 4 jobs per second. A job departing node 3 always goes to node 2.\\
\\
     (a) Calculate the mean time that a job spends in the network, from arrival to departure.\\
\\
\\
\\
     (b) You can either double the speed of server 3 or server 2. Which choice would you\\
         prefer? Justify your answer.\\
\end{quote}

\subsection*{\texttt{test2\_sample\_questions\_q95}}
\textit{test2\_sample\_questions.pdf \#95 · test2}\par\smallskip
\begin{quote}
\small\sloppy
A single server operates as follows. Arrivals occur according to a Poisson process with\\
    rate 0.15 per minute, of which 75\% are of type H and 25\% are type L. If a type L arrival\\
    occurs and the server is busy, it is lost, otherwise it begins processing. If a type H arrival\\
    occurs, there are three possibilities:\\
\\
     (i) if the server is idle, the arriving job immediately begins processing\\
     (ii) if the server is busy with a type H job, the arriving job is lost\\
    (iii) if the server is busy with a type L job, the arriving job immediately begins processing\\
          and the type L job has its processing suspended; when the type H job completes\\
          processing, the type L job resumes processing\\
\\
    The processing times for type H jobs are exponentially distributed with mean 8 minutes\\
    and the processing times of type L jobs are exponentially distributed with mean 3 minutes.\\
    Calculate the steady-state utilization of the server.\\
\end{quote}

\subsection*{\texttt{test2\_sample\_questions\_q96}}
\textit{test2\_sample\_questions.pdf \#96 · test2}\par\smallskip
\begin{quote}
\small\sloppy
A single-server system has arrivals that follow a Poisson process with rate 0.8 per minute.\\
    Processing times are exponentially distributed with mean 1 minute. There is room for\\
    three jobs in the system (arrivals that find three jobs in the system are lost). What is the\\
    throughput of this system?\\
\end{quote}

\subsection*{\texttt{test2\_sample\_questions\_q97}}
\textit{test2\_sample\_questions.pdf \#97 · test2}\par\smallskip
\begin{quote}
\small\sloppy
A closed network operates as follows. It has 4 nodes, the processing rates being 5, 10,\\
    10, 1 at the nodes, respectively. After finishing processing at node 1, jobs proceed to\\
    either node 2 or 3 with equal probability. After processing at either nodes 2 or 3, with\\
    probability 1/10 a job visits node 4, otherwise it returns to node 1. After processing\\
    at node 4, jobs return to node 1. Calculate the throughput of this system with 2 jobs\\
    circulating. If you were to improve one processing rate, which would it be? Would you\\
    expect your recommendation to change as we increase the number of jobs in the system?\\
\end{quote}

\subsection*{\texttt{test2\_sample\_questions\_q98}}
\textit{test2\_sample\_questions.pdf \#98 · test2}\par\smallskip
\begin{quote}
\small\sloppy
A continuous random variable has density\\
                                              (\\
                                                  4 3\\
                                                  15 x   1≤x≤c\\
                                    f (x) =\\
                                                  0      otherwise\\
\\
     (a) What is c?\\
     (b) Given a sample from a U [0, 1] distribution, how would one generate a sample for a\\
         random variable with density f ?\\
\end{quote}

\subsection*{\texttt{test2\_sample\_questions\_q99}}
\textit{test2\_sample\_questions.pdf \#99 · test2}\par\smallskip
\begin{quote}
\small\sloppy
Suppose that a CPU’s processing times follow the distribution\\
                                              \\
                                               0\\
                                                  x<0\\
                                     G(x) =     x2 0 ≤ x ≤ 1\\
                                              \\
                                               1  x>1\\
\\
\\
\\
     Given a sample from a U [0, 1] distribution, how can one generate a sample from the\\
     distribution G?\\
\end{quote}

\subsection*{\texttt{test2\_sample\_questions\_q100}}
\textit{test2\_sample\_questions.pdf \#100 · test2}\par\smallskip
\begin{quote}
\small\sloppy
Response times from six independent runs of a simulation are (in msec) 102.5, 101.7,\\
     103.1, 100.9, 100.5 and 102.2. What is the 95 percent confidence interval for the mean\\
     response time? How many more simulation runs would you recommend in order to reduce\\
     the width of the confidence interval by a factor of two?\\
\end{quote}

\subsection*{\texttt{test2\_sample\_questions\_q101}}
\textit{test2\_sample\_questions.pdf \#101 · test2}\par\smallskip
\begin{quote}
\small\sloppy
You have the following data: 9.50, 2.31, 6.07, 4.86, 8.91, 7.62, 4.56, 0.19, 8.21, 4.45. The\\
     confidence interval constructed for the mean was [3.94, 7.40]. What confidence level was\\
     used?\\
\end{quote}

\subsection*{\texttt{test2\_sample\_questions\_q102}}
\textit{test2\_sample\_questions.pdf \#102 · test2}\par\smallskip
\begin{quote}
\small\sloppy
You are asked to construct a 95 percent confidence interval for a mean and construct it\\
     to be (.22, 1.35). You then get another set of data, which you believe is from the same\\
     distribution, that gives an average of 1.40. Is this consistent with the confidence interval\\
     that you constructed?\\
\end{quote}

\subsection*{\texttt{test2\_sample\_questions\_q103}}
\textit{test2\_sample\_questions.pdf \#103 · test2}\par\smallskip
\begin{quote}
\small\sloppy
Calculate a 95 percent confidence interval for the following set of data: 10.95, 10.23, 10.61,\\
     10.49, 10.89, 10.76, 10.46, 10.02, 10.82, 10.44.\\
\end{quote}

\subsection*{\texttt{test2\_sample\_questions\_q104}}
\textit{test2\_sample\_questions.pdf \#104 · test2}\par\smallskip
\begin{quote}
\small\sloppy
Ten independent simulation runs give values for the average number in system of 3.36,\\
     3.65, 4.11, 3.31, 3.44, 3.50, 3.50, 3.90, 3.57, 3.50. With what level of confidence can you\\
     say that the true mean number in system is less than 4.12?\\
\end{quote}

\subsection*{\texttt{test2\_sample\_questions\_q105}}
\textit{test2\_sample\_questions.pdf \#105 · test2}\par\smallskip
\begin{quote}
\small\sloppy
An experimenter feels that for her experiment, the sample standard deviation will be no\\
     larger than 100. A 95 percent confidence interval is desired to be constructed for the\\
     mean which has a width of at most 5.0. What sample size would you recommend?\\
\end{quote}

\subsection*{\texttt{test2\_sample\_questions\_q106}}
\textit{test2\_sample\_questions.pdf \#106 · test2}\par\smallskip
\begin{quote}
\small\sloppy
Ten data points are collected, measuring the execution time of a computer job (in seconds):\\
     9.5, 2.3, 6.1, 4.9, 8.9, 7.6, 4.5, 0.2, 8.2, 4.4. With what level of confidence can we say that\\
     the mean execution time is less than 7.5 seconds?\\
\end{quote}

\subsection*{\texttt{test2\_sample\_questions\_q107}}
\textit{test2\_sample\_questions.pdf \#107 · test2}\par\smallskip
\begin{quote}
\small\sloppy
You wish to generate a sample from a distribution with density\\
                                                (\\
                                                     x/2 0 ≤ x ≤ 2\\
                                      f (x) =\\
                                                     0   otherwise\\
\\
     Using only a sample from a U [0, 1] distribution, how would you do so?\\
\end{quote}

\subsection*{\texttt{test2\_sample\_questions\_q108}}
\textit{test2\_sample\_questions.pdf \#108 · test2}\par\smallskip
\begin{quote}
\small\sloppy
Suppose that a random variable X has density\\
                                                 (\\
                                                     x   1 ≤ x ≤ b;\\
                                       f (x) =\\
                                                     0   otherwise.\\
\\
     A system has implemented a PRNG that generates samples from a U [0, 10] distribution\\
     (note that this is not U [0, 1]). Given a sample u from this PRNG, how can you generate\\
     a sample of X?\\
\end{quote}

\subsection*{\texttt{test2\_sample\_questions\_q109}}
\textit{test2\_sample\_questions.pdf \#109 · test2}\par\smallskip
\begin{quote}
\small\sloppy
Five independent replications of a simulation are run to estimate the mean response time\\
     for a server. The replications had average response times of 124, 124, 128, 130, 127. With\\
     95 percent confidence, can we say that the true mean response time is between 120 and\\
     132? Justify your answer. Would your answer change if the values were the response\\
     times of consecutive departures from the server?\\
\end{quote}

\subsection*{\texttt{test2\_sample\_questions\_q110}}
\textit{test2\_sample\_questions.pdf \#110 · test2 · T/F}\par\smallskip
\begin{quote}
\small\sloppy
True or False.\\
\\
      (a) In an M/M/k system, the expected waiting time for jobs that must wait depends\\
          only on ρ, the load on the system.\\
\end{quote}

\subsection*{\texttt{test2\_sample\_questions\_q111}}
\textit{test2\_sample\_questions.pdf \#111 · test2}\par\smallskip
\begin{quote}
\small\sloppy
(a) Given only a PRNG that returns samples from a U [0, 1] distribution, how would you\\
         generate samples from a random variable with density:\\
\\
                                             fX (x) = x2 + 2/3\\
\\
          for 0 ≤ x ≤ 1 and 0 otherwise?\\
\\
      (b) The following single-server system is simulated. There is an infinite buffer, schedul-\\
          ing is FCFS, arrivals follow a Poisson process with rate 1 per minute, and processing\\
          times are distributed according to the random variable in (a). Running the simu-\\
          lation for 10 independent replicas gives the following average response times: 0.94,\\
          1.33, 1.12, 1.41, 0.90, 1.49, 1.06, 1.49, 1.05, 1.05. Is it plausible that the simulation\\
          code is correct? Justify your answer.\\
\end{quote}

\subsection*{\texttt{test2\_sample\_questions\_q112}}
\textit{test2\_sample\_questions.pdf \#112 · test2}\par\smallskip
\begin{quote}
\small\sloppy
Consider the following network. There are three nodes, each having a single processor\\
     and exponentially distributed processing times. Departures from node 1 are equally likely\\
     to go to the other two nodes. Departures from node 2 go to node 1 with probability 0.7,\\
     otherwise they go to node 3. Departures from node 3 go to node 1 with probability 0.9,\\
     otherwise they return to node 3. The mean processing times at the three nodes are 1.5,\\
     1 and 2, respectively. Two jobs circulate in the network.\\
\\
      (a) What is the system bottleneck?\\
      (b) Calculate the throughput at node 1.\\
      (c) Calculate the expected number of jobs at node 2.\\
\end{quote}

\subsection*{\texttt{test2\_sample\_questions\_q113}}
\textit{test2\_sample\_questions.pdf \#113 · test2}\par\smallskip
\begin{quote}
\small\sloppy
A single server system operates as follows. Arrivals occur according to a Poisson process\\
     with rate λ per time unit. Processing times are exponentially distributed with mean 1\\
     time unit. Costs are as follows: 1 per job that must wait to be processed, 1 per time unit\\
     when the processor is busy, and 0.3 per time unit when the processor is idle. You have a\\
     budget of 1.5 per time unit. What is the maximum value of λ allowed?\\
\end{quote}

\subsection*{\texttt{test2\_sample\_questions\_q114}}
\textit{test2\_sample\_questions.pdf \#114 · test2}\par\smallskip
\begin{quote}
\small\sloppy
Consider the following system consisting of two servers.\\
\\
\\
\\
      (i) Type 1 arrivals follow a Poisson process with rate 2 per minute. If there is an idle\\
          server, a Type 1 arrival is assigned to it (if both servers are idle, one server is chosen\\
          randomly), otherwise it is rejected. Processing times of Type 1 jobs are exponentially\\
          distributed with mean 20 seconds.\\
      (ii) Type 2 arrivals follow a Poisson process with rate 1 per minute. Type 2 arrivals\\
           require the use of both servers simultaneously, for a processing time that is exponen-\\
           tially distributed with mean 30 seconds. If the system is busy with a Type 2 job, an\\
           arriving Type 2 job is rejected. If the system is busy with one or more Type 1 jobs,\\
           the Type 1 jobs are removed from the system and the Type 2 job begins processing.\\
\\
     Calculate the throughput of both Type 1 and Type 2 jobs.\\
\end{quote}

\subsection*{\texttt{test2\_sample\_questions\_q115}}
\textit{test2\_sample\_questions.pdf \#115 · test2}\par\smallskip
\begin{quote}
\small\sloppy
Consider the following network. Arrivals to the first node follow a Poisson process with\\
     rate 1 per time unit. After processing at node 1, jobs go to node 2 with probability 0.3\\
     and to node 3 with probability 0.6; otherwise, they leave the system. After processing at\\
     node 2 or 3, jobs return to node 1. There is a single server at each node, where processing\\
     times are exponentially distributed with rates 12 per time unit (node 1); 4 per time unit\\
     (node 2); and 8 per time unit (node 3).\\
\\
      (a) Calculate the expected number of jobs at node 3.\\
      (b) Suppose that all three servers are busy. What is the probability that an arrival\\
          occurs to node 1 (from either inside or outside of the network) before the job being\\
          processed at node 1 completes?\\
\end{quote}

\subsection*{\texttt{test2\_sample\_questions\_q116}}
\textit{test2\_sample\_questions.pdf \#116 · test2}\par\smallskip
\begin{quote}
\small\sloppy
Parts (a) and (b) are unrelated.\\
\\
      (a) Given only a PRNG that returns samples from a U [0, 1] distribution, how would you\\
          generate samples for a random variable X with distribution\\
                                           FX (x) = 1 − e−(x/3)\\
\\
      (b) You are given simulation output for an M/M/2/2 system with λ = 1 and µ = 1. The\\
          metric estimated is the mean number of jobs in the system. Ten independent runs\\
          of the simulation give the following values for average number of jobs in the system:\\
          1.026, 1.062, 0.751, 1.065, 0.953, 0.739, 0.811, 0.919, 1.083, 1.086. What are your\\
          conclusions as to whether the simulation was implemented correctly?\\
\end{quote}

\subsection*{\texttt{test2\_sample\_questions\_q117}}
\textit{test2\_sample\_questions.pdf \#117 · test2}\par\smallskip
\begin{quote}
\small\sloppy
A single server operates as follows. Arrivals occur according to a Poisson process with\\
     rate 0.8 per time unit. Processing times are exponentially distributed with mean 1 time\\
     unit. Costs are as follows: 1 per job that must wait before beginning processing, 1 per\\
     time unit when the server is busy, and 0.3 per time unit when the server is idle. Should\\
     a second server be added? (The busy/idle costs for the second server are the same as for\\
     the first server and the two servers would share a common queue.)\\
\end{quote}

